% Dies ist die Hauptdatei, von der aus das Gesamtdokument erzeugt wird.  Diese
% Datei sollten Sie zunächst umbenennen, damit sie keinen generischen Namen
% hat!  Dann wird sie mit LuaLaTeX kompiliert.

% In der Datei defs.tex werden alle globalen LaTeX-spezifischen Einstellungen
% vorgenommen
% Jede LaTeX-Datei beginnt mit der Dokumentenklasse.  Für diese Vorlage wurde
% die Klasse "scrreprt" von KOMA-Script gewählt, die in etwa der
% Standardklasse "report" entspricht, allerdings wesentlich mehr Möglichkeiten
% bietet und im gewissen "moderner" ist.  Eine sehr ausführliche Dokumentation
% zu KOMA-Script findet man unter der folgenden Adresse:
% http://mirrors.ctan.org/macros/latex/contrib/koma-script/doc/scrguide.pdf
\documentclass[
% die Schriftgröße - sollten Sie nicht ändern
    fontsize=12pt,
% das Papierformat, also DIN A4
    paper=A4,
% Literaturverzeichnis ins Inhaltsverzeichnis
    bibliography=totoc,
% andere Verzeichnisse ebenfalls ins Inhaltsverzeichnis
    listof=totoc,
% abgesetzte Formeln linksbündig
    fleqn,
% für die Satzspiegelkonstruktion - siehe KOMA-Doku
    DIV=12,
% Bindekorrektur (linker Rand) - evtl. anpassen
    BCOR=1mm,
% die im Text verwendeten Sprachen (u.a. für das Paket babel); die
% letztgenannte (!) Sprache ist die Standardsprache; "n"german steht für die
% neue Rechtschreibung
    english,
% weil (s.u.) das Paket geometr verwendet wird
    usegeometry,
% wie Absätze gesetzt werden: ohne Einzug, halbe Zeile Abstand
    parskip=half-
]{scrreprt}

\usepackage{subcaption}
% Beschriftungen für Tabellen kommen linksbündig über die Tabelle
\KOMAoption{captions}{tableheading}
\setcaptionalignment[figure]{c}
\setcaptionalignment[table]{l}

% wird für die Titelseite benötigt
\usepackage{geometry}

% Standardpaket für Lokalisation, siehe Option "ngerman" oben
\usepackage{babel}
% Laden von optimierten Trennmustern
\babelprovide[hyphenrules=ngerman-x-latest]{ngerman}

% Standardpaket für mathematische Zusatzfunktionen; wenn Sie keine
% mathematischen Formeln brauchen, können Sie diese Zeile löschen
\usepackage{amsmath}

% die Hauptschrift Libertinus
\usepackage{libertinus-otf}
% die "Schreibmaschinenschrift" Anonymous Pro, angepasst
\usepackage{AnonymousPro}
\setmonofont{AnonymousPro}[Scale=MatchLowercase,FakeStretch=0.85]

% etwas größerer Zeilenabstand als im Buchsatz
\linespread{1.1}

% Paket für Feinkorrekturen an der Typographie, das für ein ausgewogeneres
% Schriftbild sorgt
\usepackage{microtype}

% Paket für kontextsensitive Anführungszeichen
\usepackage{csquotes}
% Shortcut, damit aus dem eigentlich falschen Zeichen " richtige
% Anführungszeichen je nach Sprache werden
\MakeOuterQuote{"}

% Paket, das den Befehl \includegraphics ermöglicht
\usepackage{graphicx}

% komfortablere Aufzählungen als in Standard-LaTeX; ein Beispiel findet man in
% chap3.tex
\usepackage{enumitem}

% Paket für mehr als die üblichen Standardfarben
\usepackage[dvipsnames]{xcolor}
% Definition der "Hausfarben" der HAW
\definecolor{haw}{HTML}{003CA0}
\definecolor{haw2}{HTML}{0096D2}
\definecolor{haw3}{HTML}{A0BEDC}

% typographisch anspruchsvolle Tabellen; siehe chap3.tex
\usepackage{booktabs}

% zum Erstellen des Literaturverzeichnisses; der gängige Stil APA ist hier
% bereits eingestellt
\usepackage[style=apa]{biblatex}
% eine Beispieldatei für ein Literaturverzeichnis
\addbibresource{demo.bib}

% für die Erzeugung der Grafiken in chap3.tex; wenn Sie PGF/TikZ nicht
% verwenden wollen, können Sie diese Zeilen entfernen
\usepackage{tikz}
% Zusatzbibliotheken für TikZ, die in den genannten Beispielen verwendet
% werden
\usetikzlibrary{calc,intersections,angles,3d}

% für die Erzeugung des Codeblocks in chap3.tex; wenn in Ihrer Arbeit keine
% Codeblöcke vorkommen, können Sie diese Zeilen entfernen
\usepackage{listings}
% Anpassung des Erscheinungsbildes des Codeblocks; mehr dazu in der
% Dokumentation des Pakets "listings"
\lstdefinestyle{mystyle}{
    backgroundcolor=\color{gray!20},
    keywordstyle=\color{haw2},
    numberstyle=\footnotesize\color{haw},
    basicstyle=\ttfamily\small,
    captionpos=t,
    frame=single,
    framerule=0pt,
    keepspaces=true,
    numbers=left,
    numbersep=6pt,
    belowcaptionskip=1em,
    aboveskip=\bigskipamount,
}
\lstset{style=mystyle}
% damit es "Codeblock" und nicht "Listing" heißt
\renewcommand{\lstlistingname}{Codeblock}

% für die Verlinkung innerhalb des PDF-Dokuments, für PDF-Lesezeichen und
% PDF-Metadaten; dieses Paket sollte üblicherweise immer als letztes geladen
% werden
\usepackage[colorlinks=true,allcolors=haw,hyperfootnotes=false,pageanchor=true,linktoc=all]{hyperref}

% für die Druckversion können Sie die obige Zeile durch die folgende ersetzen,
% damit Links nicht blau dargestellt werden:
% \usepackage[draft]{hyperref}

% Metadaten des PDF-Dokumentes; setzen Sie hier Ihren eigenen Namen sowie den
% Titel Ihrer Arbeit ein
\hypersetup{pdfauthor={Prof. Dr. Edmund Weitz}}
\hypersetup{pdftitle={Handreichung zur Formatierung von Bachelorarbeiten}}


% Wenn das Kommentarzeichen entfernt wird, kann man mit einem Befehl wie
%
% \includeonly{chap2,chap3}
%
% erreichen, dass nur ausgewählte Dateien kompiliert werden.  Das ist für die
% Arbeit an umfangreichen Dokumenten hilfreich, weil es Zeit spart.  Für das
% Erstellen des fertigen Dokuments muss der Befehl natürlich wieder
% auskommentiert werden, damit alle Referenzen aktuell sind und die
% Seitenzahlen stimmen.

% Hier beginnt das eigentliche Dokument.  Sie können weitere Dateien
% hinzufügen und natürlich auch vorhandene weglassen.  Die vorhandene
% Dateistruktur ist lediglich als Beispiel gedacht.
\begin{document}
% In dieser Umgebung wird die Titelseite separat vom restlichen Text gesetzt
\begin{titlepage}
  % andere Seitenränder als im Rest der Arbeit
  \newgeometry{lmargin=2cm,tmargin=7mm,rmargin=5mm,bmargin=1cm}
  % die "Hausfarbe" der HAW; diese und die folgenden Einstellungen sind lokal
  % und gelten nur innerhalb der Umgebung "titlepage"
  \color{haw}
  % Blocksatz für die Titelseite deaktivieren
  \raggedright
  % Logo rechtsbündig setzen
  \hfill\includegraphics[width=7cm]{images/HAW_Marke_RGB_300dpi}\\

  % vertikaler Abstand
  \vspace{5cm}

  % Wahl der "Hausschrift" Open Sans der HAW, die als Schrift auf Ihrem
  % Rechner installiert sein muss
  \setmainfont{Open Sans}
  % etwas kleiner als üblich
  \small
  % fett und in Majuskeln
  \textbf{BACHELORARBEIT}

  % vertikaler Abstand
  \vspace{8mm}

  % der Titel der Arbeit als "Seite in der Seite"; natürlich müssen Sie hier
  % Ihren Titel eintragen
  \begin{minipage}{0.8\linewidth}
    % Wahle der zweiten "Hausschrift" der HAW, die ebenfalls auf Ihrem Rechner
    % bereits vorhanden sein muss
    \setmainfont{Martel Heavy}
    % ziemlich große Schrift
    \LARGE
    % [1mm] steht jeweils für einen etwas größeren Durchschuss

    Detecting Geometric Primitives\\[1mm]
    in Depth Data from the\\[1mm]
    Google ARCore Depth API\\[1mm]
    \,\rule{11mm}{1.2mm}
  \end{minipage}

  % vertikaler Abstand, überraschenderweise
  \vspace{1cm}

  % hier korrektes Datum und Ihren Namen eingeben
  vorgelegt am 01.\ Februar 1900\\
  Hugo Protsch

  % letzter vertikaler Abstand für heute
  \vspace{5cm}

  % noch eine "Seite in der Seite", etwas nach rechts geschoben
  \hspace*{37mm}
  \begin{minipage}{0.5\linewidth}
    % Namen und Titel der beiden Prüfer eintragen
    \begin{tabular}{@{}ll}
      Erstprüferin: & Prof. Dr. Philipp Jenke\\[-.3mm]
      Zweitprüfer: & Prof. Dr. Olaf Zukunft\\
    \end{tabular}\\

    % noch ein horizontaler Strich
    \,\rule{9mm}{1mm}\\[1.5mm]

    \textbf{HOCHSCHULE FÜR ANGEWANDTE}\\
    \textbf{WISSENSCHAFTEN HAMBURG}\\
    Department Informatik\\
    Berliner Tor 7\\
    20099 Hamburg
  \end{minipage}
\end{titlepage}
% setzt die Geometrie wieder auf die Standardwerte zurück
\restoregeometry

% für die Seite mit dem Abstract keine Seitenzahl ausgeben
\thispagestyle{empty}



\thispagestyle{empty}
\clearpage

% In der Titelei werden römische Ziffern für die Seitenzahlen verwendet;
% gleichzeitig wird durch diesen Befehl die aktuelle Seitenzahl auf eins
% gesetzt
\pagenumbering{Roman}

% Inhaltsverzeichnis
\tableofcontents

% Abbildungsverzeichnis, kann evtl. weggelassen werden
\listoffigures

% Tabellenverzeichnis, kann evtl. weggelassen werden
%\listoftables

% weitere Verzeichnisse (z.B. Codeblöcke) sind theoretisch möglich

% neue Seite, vorsichtshalber
\clearpage
% ab jetzt arabische Ziffern und wieder auf eins setzen
\pagenumbering{arabic}
% Inhalte in Dateien, die mit \include eingefügt werden, fangen immer auf
% einer neuen Seite an.  Dort sollte also beispielsweise wie hier ein neues
% Kapitel (\chapter) anfangen.

\chapter{Introduction}

Basic description of AR


\chapter{Die einzelnen Teile der Arbeit}

% Hier sieht man, wie mit \ref Bezug genommen wird auf Kapitel, Tabellen und
% Abschnitte.

In diesem Kapitel geht es um die wesentlichen Teile, die eine Abschlussarbeit
haben sollte.  Die Dateien, auf die hier Bezug genommen wird, werden in in
Tabelle~\ref{tab-files} vorgestellt und in Kapitel~\ref{ch-tech} näher
erläutert.  Die Vorlage wurde absichtlich in relativ viele Dateien zerlegt,
um zu demonstrieren, wie man mit dem Befehl \verb|\includeonly| nur Teile der
Arbeit kompilierern kann, um Zeit zu sparen.  Vor dem Druck muss natürlich das
komplette, aus allen Dateien bestehende Dokument kompiliert werden (ggf.\
mehrfach), damit alle Querverweise (siehe Abschnitt~\ref{sec-ref}) korrekt
sind.

% Ein Beispiel für eine Tabelle; die Tabelle selbst wird mit der
% Standardumgebung \tabular gesetzt.  Die Umgebung \table sorgt u.a. für das
% "Gleiten" und die Aufnahme in das Tabellenverzeichnis.
\begin{table}[!ht]
  % Beschriftung der Tabelle
  \caption{Dateien der Vorlage}
  % Zentrieren der Tabelle
  \centering
  \begin{tabular}{ll}
    % unterschiedliche Strichbreiten für "top", "mid" und "bottom"
    \toprule
    Dateiname & Zweck \\
    \midrule
    \texttt{VorlageBA.tex} & Hauptdatei \\
    \texttt{defs.tex} & Laden von Paketen und Setzen von Optionen \\
    \texttt{title.tex} & Titelseite und Abstract \\
    \texttt{toc.tex} & Inhaltsverzeichnis und optionale Verzeichnisse \\
    \texttt{chap1.tex} bis \texttt{chap3.tex} & erstes, zweites und drittes Kapitel \\
    \texttt{appendix.tex} & Literaturverzeichnis, Anhang \\
    & und Eigenständigkeitserklärung \\
    \texttt{demo.bib} & exemplarische Bibliographie \\[3pt]
    \verb|HAW_Marke_RGB_300dpi.jpg| & HAW-Logo für Titelseite \\[3pt]
    \texttt{bitmap.png}, \texttt{vector.pdf} & \\
    und \texttt{euler.py} & Beispieldateien \\
    \bottomrule
  \end{tabular}
  % frei wählbarer Name für \ref
  \label{tab-files}
\end{table}

\section{Titelei}

Mit dem Begriff \textit{Titelei} bezeichnet man im Buchwesen den Teil eines
Buches, der dem eigentlichen Inhalt vorangestellt ist.  In dieser Vorlage
dient der Absatz, den Sie gerade lesen, im Wesentlichen aber nur als Vorwand,
eine weitere Unterebene einzufügen.

\subsection{Titelseite}

Die Titelseite befindet sich in der Datei \texttt{title.tex} und ist der
\href{https://www.haw-hamburg.de/hochschule/hochschuleinheiten/presse-und-kommunikation/corporate-design/}{Vorlage}
der HAW nachempfunden.  Sie müssen natürlich in dieser Datei den Titel, die
Namen und das Datum anpassen.  Außerdem müssen die
\href{https://www.haw-hamburg.de/fileadmin/zentrale_PDF/Zentrale_Dokumente/Corporate_Design/HAW_Schriftenpaket.zip}{Hausschriften}
der HAW \href{https://support.microsoft.com/de-de/office/hinzuf\%C3\%BCgen-einer-schriftart-b7c5f17c-4426-4b53-967f-455339c564c1}{installiert} sein.

\subsection{Abstract}\label{sec-abstract}

Nach der Titelseite folgt der Abstract, der sich ebenfalls in der Datei
\texttt{title.tex} befindet.  Hier ersetzen Sie den Beispieltext durch eine
möglichst aussagekräftige Zusammenfassung Ihrer Arbeit.

\section{Inhaltsverzeichnis und andere Verzeichnisse}\label{sec-lists}

% Hier sieht man ein Beispiel für ein Buchzitat mit einer Kapitelangabe
Nach dem Abstract (siehe Abschnitt~\ref{sec-abstract}) folgt das
Inhaltsverzeichnis, das von \LaTeX\ automatisch erzeugt wird.  Exemplarisch
wird in dieser Vorlage auch gezeigt, wie man Abbildungs- und
Tabellenverzeichnisse generieren könnte; siehe dazu die Datei
\texttt{toc.tex}.  Das ist in der Regel aber nur dann nötig, wenn es in Ihrer
Arbeit sehr viele Abbildungen bzw.\ Tabellen gibt.  Bei Bedarf können auch
Verzeichnisse von Codeblöcken, mathematischen Formeln oder anderen Objekten
erzeugt werden; mehr dazu in \parencite[Kap.\ 12]{voss}.  Hier gilt aber
ebenfalls, dass das normalerweise nicht nötig sein sollte.
Auch ein Glossar oder ein Abkürzungsverzeichnis ist in einer Bachelorarbeit
eher unüblich.  Sprechen Sie ggf.\ mit Ihrer Erstprüferin oder Ihrem
Erstprüfer ab, was wirklich gebraucht wird.

\section{Gliederungsebenen}

Eine Bachelorarbeit sollte in der Regel mit maximal drei Gliederungsebenen
auskommen.  In der Dokumentenklasse der Vorlage entspricht das den Befehlen
\verb|\chapter|, \verb|\section| und \verb|\subsection|, die auch für eine
automatische Aufnahme der jeweiligen Abschnitte ins Inhaltsverzeichnis sorgen.
Weitere Gliederungsebenen verringern typischerweise die Lesbarkeit des
Dokuments.  Falls Sie das bei Ihrer Arbeit trotzdem für nötig halten, sprechen
Sie es vorher mit der Erstprüferin bzw.\ dem Erstprüfer ab.

Auf jeder Gliederungsebene sollte es mindestens zwei Abschnitte derselben
Hierarchie geben, da ansonsten eine Gliederung auf dieser Ebene sinnlos wäre.
Wenn Sie also einen Abschnitt 2.3.1 haben, dann muss es mindestens einen
weiteren Abschnitt 2.3.2 geben; anderenfalls sollte alles unter 2.3 stehen.

% Hier sieht man am Ende des Absatzes, dass man für Auslassungspunkte nicht
% einfach "..." tippt.  Man beachte auch die Tilde davor: das "geschützte
% Leerzeichen".
Außerdem sollten einzelne Abschnitte einen Umfang haben, der die entsprechende
Gliederung rechtfertigt.  Dieser Text, in dem Abschnitte meistens nur aus
wenigen Sätzen bestehen, ist dafür ein schlechtes Beispiel!\footnote{Darum
  wirken die Abstände und Seitenaufteilungen in dieser Vorlage auch
  vergleichsweise unruhig.  Bei längeren Texten wird es besser aussehen.} Es
handelt sich allerdings auch um eine Vorlage und nicht um eine
Bachelorarbeit~\dots

In der Vorlage ist jedes Kapitel in eine eigene Datei ausgelagert.  Beachten
Sie, dass der Befehl \verb|\include| grundsätzlich eine neue Seite anfängt.
Unterabschnitte von Kapiteln müssen daher mit \verb|\input| eingefügt werden,
wenn sie in separaten Dateien liegen sollen.

\section{Literaturverzeichnis}

% Beachten Sie, dass bei Abkürzungen wie "z.B." (siehe unten) oder "u.a." ein
% "kleines" Leerzeichen hinter den ersten Punkt gehört
Die Vorlage enthält ein exemplarisches Literaturverzeichnis, das nach dem
sogenannten APA-Standard formatiert ist und das als Beispiele einige Bücher,
einen Fachartikel und eine Internetquelle umfasst.  Das Literaturverzeichnis
befindet sich am Ende der Arbeit vor dem Anhang.  Falls Ihre Erstprüferin oder
Ihr Erstprüfer ein anderes Format wünscht, ist das mit dem verwendeten Paket
\href{https://www.ctan.org/pkg/biblatex}{Bib\LaTeX} ohne große Probleme zu
realisieren, siehe z.\,B. \parencite[Kap.\ 13]{voss}.  Wichtig ist, dass das
Literaturverzeichnis vollständig ist und Sie nur die Quellen angeben, die Sie
im Rahmen Ihrer Bachelorarbeit wörtlich zitiert oder sinngemäß wiedergegeben
haben.  Literatur, die Sie lediglich zur Vorbereitung genutzt haben, gehört
nicht in das Literaturverzeichnis.\footnote{Wenn Sie die Voreinstellungen der
  Vorlage verwenden, werden aber ohnehin nur die Quellen ins
  Literaturverzeichnis übernommen, die auch zitiert werden.}

Wenn Sie die von Ihnen verwendeten Texte nach dem Muster der Datei
\texttt{demo.bib} eingeben, wird das Literaturverzeichnis automatisch
einheitlich dargestellt und alphabetisch sortiert.  Die Literaturverwaltung
\href{https://de.wikipedia.org/wiki/Citavi}{\textsc{Citavi}}, für die es eine
Hochschullizenz gibt, kann Daten im sogenannten Bib\TeX-Format ausgeben.  (Das
ist das in \texttt{demo.bib} verwendete Format.)

Verwenden Sie wissenschaftliche Quellen.  (Hinweis: Wikipedia gilt nicht als
wissenschaftliche Quelle.)  Wenn die Angabe von Internetquellen unumgänglich
ist, muss das Datum des letzten Aufrufs angegeben werden.  Ein Beispiel dafür
finden Sie in der Vorlage.

Wie man zitiert, wird in Abschnitt~\ref{sec-litref} gezeigt.

\section{Anhang}\label{sec-app}

Falls zu Ihrer Arbeit Datenreihen, Quelltexte, transkribierte Interviews oder
weitere ergänzende Informationen gehören, dann gehören diese in einen Anhang
ganz am Ende.  Ob ein Anhang notwendig ist und welchen Umfang er haben sollte,
sollten Sie vorab mit Ihrer Erstprüferin oder Ihrem Erstprüfer absprechen.
Häufig ist es sinnvoller, die Daten auf einem Datenträger der Arbeit
beizulegen.

Die Vorlage enthält einen exemplarischen Anhang in der Datei
\texttt{appendix.tex}.

\section{Eigenständigkeitserklärung}

Die letzte Seite der Arbeit ist die Eigenständigkeitserklärung, die Sie in der
Datei \texttt{appendix.tex} finden.  Tragen Sie hier den Titel Ihrer Arbeit
sowie den Ort und das Datum ein.  Alle gedruckten Exemplare werden dann
oberhalb des Datums von Ihnen unterschrieben.

\chapter{TBD}

% -*- coding: utf-8 -*-

% Ausgabe des Literaturverzeichnisses; ohne weitere Optionen werden nur die
% Bücher und Artikel ausgegeben, die in der Arbeit auch zitiert werden.
\printbibliography

% markiert den Anfang des Anhangs
%\appendix

% ein Kapitel, das nicht numeriert, aber trotzdem ins Inhaltsverzeichnis
% aufgenommen wird
%\addchap{Anhang}
%
%Hier beginnt der Anhang.  Siehe die Anmerkungen zur Sinnhaftigkeit eines
%Anhangs in Abschnitt~\ref{sec-app} auf Seite~\pageref{sec-app}.
%
%Der Anhang kann wie das eigentliche Dokument in Kapitel und Abschnitte
%unterteilt werden.  Der Befehl \verb|\appendix| sorgt im Wesentlichen nur für
%eine andere Nummerierung.

% neue Seite
\clearpage

% keine Seitenzahl
\thispagestyle{empty}

% keine Nummerierung, keine Aufnahme ins Inhaltsverzeichnis
%\section*{Eigenständigkeitserklärung}

% Hier müssen Sie natürlich den Titel der Arbeit sowie Ort und Datum ersetzen:
%Hiermit versichere ich, dass ich die vorliegende Bachelorarbeit mit dem Titel
%\begin{center}
%  \textbf{Viele zufällige Zahlen}
%\end{center}
%selbstständig und nur mit den angegebenen Hilfsmitteln verfasst habe.  Alle
%Passagen, die ich wörtlich aus der Literatur oder aus anderen Quellen wie
%z.\,B. Internetseiten übernommen habe, habe ich deutlich als Zitat mit Angabe
%der Quelle kenntlich gemacht.
%
%\vspace{2cm}
%
%Hamburg, 21.\ Dezember 1940

\end{document}